\section{Conclusion}
\label{sec:conclusion}

This study examined the number of candidate features per split as a design variable of hybrid Wi-Fi RTT--RSS Random Forest localization, where prior work has used the full-feature setting, i.e., bagging of regression trees. On three public environments, restricting each split to $70\%$ of the RTT--RSS features reduced the error of the reproduced full-feature baseline by $14.3\%$, $7.1\%$, and $5.1\%$ on the official RP-disjoint test sets, with RP-clustered bootstrap intervals that exclude zero for Building and Office and a lower error in all ten seeds in every environment, and by $12$--$22\%$ under RP-grouped cross-validation.

The ratio analysis explains these numbers. Full-feature bagging was never the best ratio. The optimal ratio decreased in all three environments as the training radio map was thinned, and the gain of subsampling reached $25.5\%$, which is consistent with per-split randomization acting as regularization for the harder interpolation problem posed by sparse maps; extremely randomized trees, which are already strongly randomized, gain at most $0.8\%$ from subsampling. The same effect biases RP-grouped cross-validation toward small ratios, while out-of-bag and sample-level validation are dominated by repeated scans of the same RP. The fixed default $\rho=0.7$ stayed within $4.1\%$ of the best ratio in every environment, compared with $19.7\%$ for bagging, and ratios between $0.5$ and $0.7$ are recommended for densely surveyed maps.

The improvement transferred. On the independent EETAC dataset, the frozen configuration was better than the baseline in 74 of 80 seed--protocol comparisons, including three of the four cross-day directions; the fourth is a tie. This evaluation required correcting a mirrored coordinate frame between the two acquisition days. Without the correction, cross-day errors exceed $5$~m; with it, they are $1.2$--$1.7$~m. The accuracy gained by subsampling also allows smaller models: 50--100-tree forests remained more accurate than the 300-tree full-feature forest while being $2.5$--$5.5$ times smaller. Accuracy remains below that of a recent graph-temporal network on the largest environment, at a small fraction of its training cost.

For uncertainty, covariance-aware ellipses describe anisotropic errors more compactly than circles, but the main lesson concerns calibration: because the scans of an RP are near-duplicates, the effective calibration sample size is the number of RPs. With 21 calibration RPs, scan-level split-conformal regions fell below the nominal coverage on average, whereas group calibration with one scan per RP attained it with $22$--$30\%$ larger regions; with 7 RPs, no valid finite region exists. Calibration should therefore be planned in terms of reference points rather than scans, with at least $1/\alpha-1$ calibration RPs.

Future work will develop a density-aware rule for the subspace ratio, cross-fitted group-conformal calibration that does not withhold RPs from the point predictor, longer-term temporal evaluations, additional buildings and devices, and on-device measurements of compact forests.
